% --------------------------------------------------------------------------
% Draft-only helper.
% Delete \AppendixPlan and replace each scaffold with final prose before
% submission.
% --------------------------------------------------------------------------
\newcommand{\AppendixPlan}[3]{%
\paragraph{Question addressed.}
#1

\paragraph{Content to include.}
#2

\paragraph{Required evidence.}
#3
}


% ==========================================================================
% Appendix A
% ==========================================================================
\section{Detailed Design and Operational Semantics}
\label{app:design-details}


% --------------------------------------------------------------------------
\subsection{Common Request State and Notation}
\label{sec:appendix:state-model}

RAVEL separates \emph{placement ownership} from engine execution state.
Before successful engine handoff, a request is Router-owned: it has one
tentative replica, contributes prospective demand to that replica, and may be
replanned subject to the mobility rules below.
Successful engine handoff transfers ownership to the serving engine and makes
the geographic placement final.
Engine-side execution may subsequently move through Prefill, Decode,
completion, or failure states, but these events do not reopen geographic
placement.

For request $i$, $a_i$ denotes its Router arrival time, $o_i$ its client
region, and $r_i(t)$ its tentative replica while Router-owned.
The version counter $v_i$ protects versioned ownership and accounting
transfers, while $h_i^{\mathrm{att}}$ identifies the current engine attempt.
Attempt identifiers isolate asynchronous completion, cancellation, and cleanup
events across retries.

The Router-side transport path distinguishes
\texttt{WAITING\_MOVABLE} from the transient \texttt{COMMITTING} state around
engine handoff.
A dispatch claim enters \texttt{COMMITTING}; this is an ownership-transition
attempt rather than a second placement decision.
If target registration or pre-handoff posting setup fails, Router ownership is
restored or cleaned without reporting a committed placement.
Placement becomes final only after the target accepts registration and the
engine-facing posting task is created and tracked successfully.
RAVEL does not perform ordinary geographic rebinds after this boundary.

These operational rules refine the three invariant classes from
Section~\ref{sec:design:overview}.
Specifically, RAVEL ensures that
(i) each Router-owned request has exactly one tentative owner;
(ii) its prospective demand is represented exactly once and follows
reassignment;
(iii) only Router-owned requests may change geographic placement;
(iv) replanning is causal and bounded in observation time and cohort size;
and (v) only successful engine handoff commits placement.
Attempt identifiers additionally isolate asynchronous callbacks across
retries: a callback from attempt $h$ cannot modify state belonging to a newer
attempt $h'>h$.


% --------------------------------------------------------------------------
\subsection{Replica SLO Quote Construction}
\label{sec:appendix:quotes}

A placement decision consumes the causal quote
\begin{equation}
    Q_{i,r}(t)
    =
    \left(
        \Phi_{i,r}(t),
        \widehat{T}^{\mathrm{obj}}_{i,r}(t),
        \mathbf{z}_{i,r}(t)
    \right),
    \label{eq:appendix:quote-tuple}
\end{equation}
where $\Phi_{i,r}(t)$ is predicted SLO feasibility,
$\widehat{T}^{\mathrm{obj}}_{i,r}(t)$ is the class-specific ordering objective,
and $\mathbf{z}_{i,r}(t)$ contains secondary capacity-pressure and stability
signals.
Quotes use only information available at decision time and are not formal SLO
guarantees.
Their inputs include request-visible SLO metadata, directed WAN measurements,
calibrated per-cluster service profiles, observed engine state, Router-owned
queue state, and tentative assignments already selected in the current
planning pass.
The evaluated request's realized output length is never an input to its quote.

\paragraph{Calibrated service curves.}
For cluster $c$ and active Decode concurrency $n$, the service profile stores
an affine Prefill model
\begin{equation}
    P_{c,n}(p)=b_{c,n}+a_{c,n}p
    \label{eq:appendix:prefill-profile}
\end{equation}
and a measured post-first-token Decode time per token $d_c(n)$.
Profiles are calibrated independently for each serving cluster and carry an
engine fingerprint.
At runtime, RAVEL linearly interpolates between measured concurrency points
and clamps values outside the calibrated range to the nearest endpoint.
Because $P_{c,n}$ is measured while $n$ Decode sequences are active, RAVEL
does not add a separate ad-hoc Decode-interference penalty to the Prefill
estimate.

\paragraph{Queue-ahead Prefill expansion.}
Let $p_i$ be the full prompt-token charge of request $i$.
Without engine-certified cached-token residency, shadow prefix locality is used
only as a late placement tie-breaker and does not reduce $p_i$.

For candidate replica $r$, let $E_r$ denote observed engine-pending Prefill
tokens, $B_{i,r}$ the prompt tokens in Router-owned requests ordered no later
than $i$, $N_r^{\mathrm{eng}}$ the number of engine-pending requests, and
$N_{i,r}^{\mathrm{q}}$ the corresponding number of Router-owned queue-ahead
requests.
With $n_r$ currently running Decode sequences, RAVEL estimates
\begin{equation}
    \widehat{S}^{\mathrm{pre}}_{i,r}
    =
    a_{c(r),n_r}
    \left(E_r+B_{i,r}+p_i\right)
    +
    b_{c(r),n_r}
    \left(N_r^{\mathrm{eng}}+N_{i,r}^{\mathrm{q}}+1\right).
    \label{eq:appendix:prefill-service}
\end{equation}
During cohort replanning, $B_{i,r}$ and $N_{i,r}^{\mathrm{q}}$ additionally
include tentative assignments selected earlier in the same planning pass.
Thus, later candidates observe prospective work induced by earlier tentative
choices rather than a stale snapshot.

Using the notation of Section~\ref{sec:design:feasibility}, the TTFT estimate
is
\begin{equation}
    \widehat{\mathrm{TTFT}}_{i,r}(t)
    =
    (t-a_i)
    +
    \rho_{o_i,c(r)}
    +
    \widehat{S}^{\mathrm{pre}}_{i,r},
    \label{eq:appendix:ttft}
\end{equation}
where $\rho_{o_i,c(r)}$ is the directed client-to-cluster WAN delay.
The elapsed term explicitly charges Router observation time against the
request's remaining SLO budget.

\paragraph{Prospective Decode occupancy.}
Let $\mathcal{O}_r$ be the union of request identifiers visible in the
Router-owned queue, engine-pending set, and engine-running set at replica $r$.
RAVEL forms
\begin{equation}
    m_{i,r}
    =
    \max\!\left(
        1,\;
        \left|\mathcal{O}_r\cup\{i\}\right|,\;
        N_r^{\mathrm{q}}+N_r^{\mathrm{eng}}+N_r^{\mathrm{run}}
        +\mathbf{1}[i\notin\mathcal{O}_r]
    \right),
    \label{eq:appendix:prospective-seqs}
\end{equation}
and queries the calibrated Decode curve at $m_{i,r}$, clamping at the measured
profile endpoint.
The set term prevents the same request from being charged twice when lifecycle
views overlap.

\paragraph{LATENCY feasibility.}
For a LATENCY request with TTFT and TBT limits
$D_i^{\mathrm{TTFT}}$ and $D_i^{\mathrm{TBT}}$,
\begin{equation}
    \Phi^{\mathrm{lat}}_{i,r}(t)
    =
    \mathbf{1}\!\left[
        \widehat{\mathrm{TTFT}}_{i,r}(t)
        \le D_i^{\mathrm{TTFT}}
        \;\land\;
        d_{c(r)}(m_{i,r})
        \le D_i^{\mathrm{TBT}}
    \right].
    \label{eq:appendix:latency-feasible}
\end{equation}
The two constraints are evaluated independently: a low predicted TTFT cannot
compensate for a predicted TBT violation.

\paragraph{Completion feasibility and output demand.}
For completion-oriented requests, RAVEL distinguishes request-visible expected
output demand $L_i^{\mathrm{exp}}$ from conservative risk demand
$L_i^{\mathrm{risk}}$.
The corresponding completion estimates are
\begin{align}
    \widehat{T}^{\mathrm{exp}}_{i,r}(t)
    &=
    \widehat{\mathrm{TTFT}}_{i,r}(t)
    +
    \max(0,L_i^{\mathrm{exp}}-1)d_{c(r)}(m_{i,r}),
    \label{eq:appendix:completion-exp}\\
    \widehat{T}^{\mathrm{risk}}_{i,r}(t)
    &=
    \widehat{\mathrm{TTFT}}_{i,r}(t)
    +
    \max(0,L_i^{\mathrm{risk}}-1)d_{c(r)}(m_{i,r}).
    \label{eq:appendix:completion-risk}
\end{align}
Completion feasibility uses
\begin{equation}
    \Phi^{\mathrm{comp}}_{i,r}(t)
    =
    \mathbf{1}\!\left[
        \widehat{T}^{\mathrm{risk}}_{i,r}(t)
        \le D_i^{\mathrm{comp}}
    \right],
    \label{eq:appendix:completion-feasible}
\end{equation}
while $\widehat{T}^{\mathrm{exp}}_{i,r}$ represents expected shared service
pressure.

Ordinary feasibility always uses the conservative
$L_i^{\mathrm{risk}}$ estimate.
Under the completion-pressure path in
Section~\ref{sec:appendix:completion-pressure}, a held-out semantic profile may
replace the placement-time estimate of $L_i^{\mathrm{exp}}$ using only
request-visible workflow metadata.
It does not replace $L_i^{\mathrm{risk}}$, reveal future arrivals, or use the
evaluated request's realized output length.

\paragraph{Quote procedure.}
For completeness, a replica quote performs the following operations:
\begin{quote}
\textsc{QuoteReplica}$(i,r,t)$:
(1) read calibrated service curves and currently observed engine state;
(2) charge full prompt work and Router-owned work ahead of $i$;
(3) add tentative work selected earlier in the current planning pass;
(4) compute Eq.~\ref{eq:appendix:ttft};
(5) form the request-ID-deduplicated prospective Decode count in
Eq.~\ref{eq:appendix:prospective-seqs};
(6) evaluate the applicable LATENCY or completion feasibility predicate; and
(7) return feasibility, the class-specific objective, and secondary
pressure/locality signals.
\end{quote}


% --------------------------------------------------------------------------
\subsection{Prospective Accounting Lifecycle}
\label{sec:appendix:accounting}

RAVEL uses prospective state for two purposes: to make newly arrived work
visible before commitment and to ensure that sequential decisions within one
planning pass account for tentative work selected earlier in that pass.

\paragraph{Queue and service debt.}
The queue-ahead terms in Eq.~\ref{eq:appendix:prefill-service} derive from
Router-owned and engine-visible state.
The RAVEL planner additionally maintains virtual ledgers so that newly
tentative work affects subsequent placement decisions before engine lifecycle
feedback arrives.

For a LATENCY request,
\begin{equation}
    u_{i,r}
    =
    \frac{d_{c(r)}(1)}
         {D_i^{\mathrm{TBT}}},
    \qquad
    A_r
    =
    \sum_{j:\,\mathrm{latency\ owner}(j)=r}u_{j,r}.
    \label{eq:appendix:latency-ledger}
\end{equation}
Both numerator and denominator are seconds, so $u_{i,r}$ is dimensionless and
represents one-token-period service demand.
This ledger is a placement-balancing signal rather than an estimate of the
instantaneous vLLM Decode queue.

For requests with active service accounting, define
\begin{equation}
    s^{\mathrm{own}}_{i,r}
    =
    \widehat{S}^{\mathrm{own,pre}}_{i,r}
    +
    \widehat{S}^{\mathrm{own,dec}}_{i,r},
    \qquad
    L_r
    =
    \sum_{j:\,\mathrm{service\ owner}(j)=r}
       s^{\mathrm{own}}_{j,r}.
    \label{eq:appendix:service-ledger}
\end{equation}
The Prefill component contains the request's own profiled service, including
the profile intercept, rather than queue-ahead debt.
The Decode component uses the request-visible expected service represented by
the corresponding quote.
Thus $L_r$ is measured in seconds of profiled service.

\paragraph{Exactly-once registration and transfer.}
Each persistent virtual charge is associated with an explicit request owner.
Registration is idempotent: if the corresponding owner field already exists,
the request is not added again.

For any transferable prospective ledger $Z$, reassignment from $r$ to $r'$
performs
\begin{equation}
    Z_r \leftarrow Z_r-z_{i,r},
    \qquad
    Z_{r'} \leftarrow Z_{r'}+z_{i,r'},
    \label{eq:appendix:account-transfer}
\end{equation}
where the destination charge is recomputed for $r'$.
The remove--recompute--install sequence executes atomically on the serialized
Router control path.

During cohort replanning, RAVEL first removes mobile candidates' old
transferable charges from the fixed planning base.
Candidates are then processed in planning order; after each tentative target
is selected, its new prompt, count, and service charges are added immediately
to pass-local state.
Persistent ownership is updated only after the corresponding version and
attempt checks succeed.

For completion-oriented accounting, the active service charge is released by
the matching terminal callback.
The LATENCY token-period ledger remains a placement-balancing signal; actual
Decode feasibility uses the request-ID-deduplicated prospective count in
Eq.~\ref{eq:appendix:prospective-seqs}.

\begin{table}[t]
    \centering
    \small
    \caption{Prospective-accounting lifecycle.}
    \label{tab:appendix:accounting-lifecycle}
    \setlength{\tabcolsep}{3.5pt}

    \begin{tabular}{
        @{}
        p{0.29\columnwidth}
        p{0.22\columnwidth}
        p{0.29\columnwidth}
        p{0.13\columnwidth}
        @{}
    }
        \toprule
        \textbf{Event} &
        \textbf{Router owner} &
        \textbf{Prospective} &
        \textbf{Active} \\
        \midrule
        Arrival            & Install & Install          & --      \\
        Reassignment       & Move    & Recompute + move & --      \\
        Rejected handoff   & Restore & Restore          & --      \\
        Successful handoff & Remove  & Remove           & Install \\
        Terminal event     & --      & --               & Release \\
        \bottomrule
    \end{tabular}
\end{table}
Operationally,
\begin{quote}
\textsc{TransferOrRefresh}$(i,r\!\rightarrow\!r')$:
(1) verify that $i$ is still Router-owned and its version/attempt is current;
(2) remove its old transferable charges from the planning base;
(3) recompute $Q_{i,r'}$ against fixed and pass-local state;
(4) update pass-local ledgers immediately if $r'$ is selected;
(5) atomically replace persistent owner and charge fields; and
(6) record the updated tentative quote.
A failed version or attempt check leaves the current owner unchanged.
\end{quote}


% --------------------------------------------------------------------------
\subsection{Bounded Pre-Handoff Replanning}
\label{sec:appendix:replanning}

RAVEL replans from causally observed state rather than polling for an
unbounded future window.
Scheduling operations are serialized by one asynchronous scheduler lock, and
multiple requests for a planning pass may be coalesced.

\paragraph{Replanning triggers and eligibility.}
A planning pass may be requested by:
(i) a new arrival,
(ii) mobility-hold expiration,
(iii) a relevant engine-state update,
(iv) request termination,
(v) rejected handoff, or
(vi) a completion-pressure state change.
Only Router-owned requests remain eligible for ordinary geographic
reassignment.
A candidate is excluded once it has crossed the handoff boundary, has no
alternative SLO-feasible target, or has exhausted the configured rebind limit.

\paragraph{Slack-bounded observation.}
Let $\delta_i$ be the most recently observed interarrival gap and
$\mathcal{D}_i$ the set of positive RTTs from request $i$'s origin to remote
serving clusters.
The causal burst predicate is
\begin{equation}
    \mathrm{Burst}_i
    =
    \mathbf{1}\!\left[
        \mathcal{D}_i\neq\emptyset
        \land
        \delta_i\le\max\mathcal{D}_i
    \right].
    \label{eq:appendix:burst-predicate}
\end{equation}

Let $\Delta_i$ denote the applicable planning budget:
$D_i^{\mathrm{TTFT}}$ for LATENCY requests and
$D_i^{\mathrm{comp}}$ for completion-oriented requests.
Arrival-time planning slack is
\begin{equation}
    \sigma_i^{\mathrm{arr}}
    =
    \max\!\left\{
        0,\,
        \Delta_i-\min_r
        \widehat{T}^{\mathrm{obj}}_{i,r}(a_i)
    \right\}.
    \label{eq:appendix:arrival-slack}
\end{equation}

Let $\rho_i^{\min}=\min\mathcal{D}_i$ when
$\mathcal{D}_i\neq\emptyset$.
Hold eligibility requires
(i) a non-LATENCY request,
(ii) $\mathrm{Burst}_i=1$, and
(iii) at least two SLO-feasible serving regions.
The mobility hold is
\begin{equation}
    h_i
    =
    \begin{cases}
        \min\!\left\{
            \rho_i^{\min},
            \sigma_i^{\mathrm{arr}}
        \right\},
        & \text{if $i$ is hold-eligible},\\
        0,
        & \text{otherwise}.
    \end{cases}
    \label{eq:appendix:mobility-hold}
\end{equation}
LATENCY requests are not deliberately held.
The hold changes dispatch eligibility rather than accounting: a held request
remains tentatively placed, prospectively accounted, and revisable throughout
the observation interval.

\paragraph{Movable cohort.}
Movable requests are ordered lexicographically by
\begin{equation}
    \kappa_i
    =
    \left(
        a_i+\Delta_i,\;
        a_i,\;
        id_i
    \right),
    \label{eq:appendix:planning-order}
\end{equation}
which gives a deterministic deadline-first order.
The ordinary movable-candidate bound is $K$; under completion pressure,
Section~\ref{sec:appendix:completion-pressure}, the effective bound may expand
to aggregate sequence capacity across replicas.

\paragraph{Fixed and pass-local context.}
A planning pass freezes noncandidate engine and Router context.
Fixed Router-owned requests that precede candidate $i$ under
Eq.~\ref{eq:appendix:planning-order} contribute their prompt work and request
intercept to $i$'s quote.
After a target is selected for a movable candidate, its prompt work, request
count, service charge, and applicable TBT charge are immediately added to
pass-local state before the next candidate is evaluated.

The full planning procedure is:
\begin{quote}
\textsc{Replan}:
(1) collect the bounded movable cohort and sort it by
Eq.~\ref{eq:appendix:planning-order};
(2) freeze current noncandidate engine and Router state;
(3) remove candidates' old transferable charges from the fixed planning base;
(4) for each candidate, quote every legal replica, choose the minimum
lexicographic placement key, and update pass-local state immediately;
(5) apply only still-current assignments, respecting attempt/version and
rebind guards; and
(6) dispatch newly eligible work through the ordinary dispatch frontier.
\end{quote}

Let $N$ be the movable cohort size, $R$ the number of replicas, and $F$ the
number of fixed Router-owned requests scanned to construct queue-ahead
prefixes.
The planning cost is
\begin{equation}
    O(N\log N+F\log F+NR+F),
    \label{eq:appendix:planner-complexity}
\end{equation}
plus constant-size service-profile interpolation.
The $NR$ term is the request--replica quote-comparison cost; with
$N\le K_{\mathrm{eff}}$, it corresponds to the
$O(K_{\mathrm{eff}}R)$ comparison bound stated in
Section~\ref{sec:design:replanning}.
The remaining terms account for deterministic cohort ordering and fixed-context
construction.
The bound is independent of generated token count; RAVEL does not require a
Router scheduling event for every generated token.


% --------------------------------------------------------------------------
\subsection{Deterministic Placement Key}
\label{sec:appendix:placement-key}

RAVEL avoids combining heterogeneous SLO and load quantities in a fitted
weighted sum.
The planner compares candidate replica $r$ for request $i$ using the
lexicographic key
\begin{equation}
\small
\begin{split}
    K_{i,r} =
    (&
      \neg\Phi_{i,r},\;
      O_{i,r},\;
      \ell^{\mathrm{obj}}_{i,r},\;
      \ell^{\mathrm{TTFT}}_{i,r},\\
     &A_r+u_{i,r},\;
      L_r+s^{\mathrm{own}}_{i,r},\;
      \widehat{T}^{\mathrm{obj}}_{i,r}/D_i,\;
      M_{i,r},\;
      \mathrm{moved}_{i,r},\;
      id_r
    ),
    \label{eq:appendix:placement-key}
\end{split}
\end{equation}
with lower values preferred.
Terms inapplicable to a request class use their neutral value rather than
inducing a separate class-specific placement algorithm.

The exact order can be interpreted in three groups.
The \emph{SLO-safety} prefix first excludes predicted-infeasible targets and
applies the pending-capacity overflow guard $O_{i,r}$ together with normalized
SLO lateness.
\emph{Capacity-pressure} terms then compare prospective LATENCY/TBT and
completion-service pressure.
Finally, \emph{efficiency and stability} use the predicted objective,
marginal-service cost, movement, and deterministic replica identity.
Lower-priority terms cannot override an earlier safety decision.

The normalized lateness terms are
\begin{align}
    \ell^{\mathrm{obj}}_{i,r}
    &=
    \frac{
        \max(0,\widehat{T}^{\mathrm{obj}}_{i,r}-D_i)
    }{\max(D_i,\varepsilon)},
    \label{eq:appendix:objective-lateness}\\
    \ell^{\mathrm{TTFT}}_{i,r}
    &=
    \frac{
        \max(0,\widehat{\mathrm{TTFT}}_{i,r}
                  -D_i^{\mathrm{TTFT}})
    }{\max(D_i^{\mathrm{TTFT}},\varepsilon)}.
    \label{eq:appendix:ttft-lateness}
\end{align}
$O_{i,r}$ is the pending-capacity overflow cost when that guard is applicable.

The convex marginal-service term is
\begin{equation}
    M_{i,r}
    =
    \frac{
       (W_r+s^{\mathrm{own}}_{i,r})^2-W_r^2
    }{D_i^2},
    \label{eq:appendix:marginal-work}
\end{equation}
where $W_r$ is the profiled service already represented at replica $r$ in the
current planning state.
Equation~\ref{eq:appendix:marginal-work} is the discrete increment of a
quadratic congestion potential rather than a learned weight.

Initial single-request placement uses the same SLO-safety and capacity
ordering.
Shadow prefix locality, when consulted, is restricted to a late tie-break path
and never precedes predicted feasibility or normalized SLO violation.

During cohort replanning, the stronger invariant is pass-local visibility:
after each tentative assignment, the planner updates prospective state before
quoting the next candidate.
Later candidates therefore observe demand induced by earlier choices in the
same pass rather than a stale pre-pass snapshot.


% --------------------------------------------------------------------------
\subsection{Completion-Pressure Adaptation}
\label{sec:appendix:completion-pressure}

The completion-pressure path is enabled by request-visible workflow semantics,
not by a workload or dataset name.
A request is a workflow-completion request when it is COLLECTIVE and contains
more than one stage.
RAVEL tracks the number of such arrivals, their first and latest arrival
times, the minimum observed completion SLO, whether nonworkflow traffic has
appeared, and the amount of outstanding workflow work.

Let $C_{\mathrm{seq}}$ be aggregate sequence capacity across replicas and
$D_{\min}$ the minimum completion SLO among observed workflow requests.
After at least $N_{\min}$ samples, RAVEL computes
\begin{equation}
    q_{\mathrm{obs}}
    =
    \frac{N_{\mathrm{arr}}-1}
         {t_{\mathrm{last}}-t_{\mathrm{first}}},
    \qquad
    q_{\mathrm{act}}
    =
    \frac{C_{\mathrm{seq}}}
         {\max(1\,\mathrm{s},\,0.5D_{\min})},
    \label{eq:appendix:completion-pressure-qps}
\end{equation}
where
\begin{equation}
    N_{\min}
    =
    \max\!\left(
       4,\;
       \min\!\left(
          16,\,
          \left\lfloor
             \frac{C_{\mathrm{seq,replica}}}{4}
          \right\rfloor
       \right)
    \right).
    \label{eq:appendix:pressure-min-samples}
\end{equation}

Define $\mathrm{WorkflowOnly}(t)$ to indicate that workflow-completion traffic
has been observed and no nonworkflow request has yet appeared.
The effective pressure path is active when
\begin{equation}
\mathrm{CompletionPressure}(t)
=
\mathbf{1}\!\left[
\begin{aligned}
    &N_{\mathrm{arr}} \ge N_{\min} \\
    &{}\land \mathrm{ObservedPressure}(t) \\
    &{}\land N_{\mathrm{outstanding}} > 0 \\
    &{}\land \mathrm{WorkflowOnly}(t)
\end{aligned}
\right],
\label{eq:appendix:completion-pressure}
\end{equation}
where $\mathrm{ObservedPressure}$ becomes true once
$q_{\mathrm{obs}}\ge q_{\mathrm{act}}$.
This observed-pressure bit is causal and sticky within the run.
The effective path becomes inactive when its eligibility conditions no longer
hold and is cleared after outstanding work drains to zero.

Completion pressure changes only two planner inputs:
\begin{enumerate}
    \item the movable-candidate ceiling expands from the ordinary bound $K$ to
          aggregate deployment sequence capacity; and
    \item when a held-out semantic profile covers the request context,
          placement-time $L_i^{\mathrm{exp}}$ uses the profile's expected
          output estimate.
\end{enumerate}
It does not change $L_i^{\mathrm{risk}}$, the feasibility predicate, the
placement algorithm, the dispatch path, or the commitment boundary.

\paragraph{Semantic output profile.}
The held-out profile indexes
$(\textit{stage\_id},\textit{stage\_num})$ and stores an expected output
estimate for each covered context.
Lookup follows
\begin{equation}
\begin{aligned}
(\textit{stage\_id},\textit{stage\_num})
   &\rightarrow
   \textit{stage\_id aggregate}\\
   &\rightarrow
   \text{declared routing hint}.
\end{aligned}
\label{eq:appendix:output-profile-fallback}
\end{equation}
Neither future arrivals nor the current request's realized output length
enters the lookup.
When an active completion request terminates, RAVEL releases its active
service charge and requests another serialized planning pass.


% --------------------------------------------------------------------------
\subsection{Shallow Reversible Dispatch Frontier}
\label{sec:appendix:frontier}

The RAVEL dispatch frontier limits how much Router-visible prompt work can
cross toward engine state at once.
Let
\begin{equation}
    H
    =
    \max_{c\neq c'}
       \rho_{c,c'}
    \label{eq:appendix:control-horizon}
\end{equation}
be the largest configured directed inter-cluster control RTT, and let
$B_{\mathrm{iter}}$ be the engine's configured
\texttt{max\_num\_batched\_tokens}.
For replica $r$, let $a_{c(r),0}$ be the calibrated fastest Prefill seconds per
token.
RAVEL sets
\begin{equation}
    F_r
    =
    B_{\mathrm{iter}}
    +
    \left\lceil
       \frac{H}{a_{c(r),0}}
    \right\rceil .
    \label{eq:appendix:frontier-depth}
\end{equation}
Thus, the frontier covers approximately one complete engine scheduling
quantum plus the prompt work that the fastest calibrated Prefill path could
consume during one worst-case control RTT.

At dispatch time, if $E_r$ is observed engine-pending prompt work, frontier
headroom is
\begin{equation}
    G_r
    =
    \max(0,F_r-E_r).
    \label{eq:appendix:frontier-headroom}
\end{equation}
The effective dispatch budget is the minimum of $G_r$ and the replica's
ordinary engine budget.
A Router-owned request may cross only if its full conservative prompt charge
fits the budget.
Eligible requests are ordered by applicable absolute SLO deadline, then
arrival time and request identifier.

In our five-replica deployment, the maximum directed RTT is
$H=37.9$\,ms.
We use $B_{\mathrm{iter}}=16384$ tokens,
\texttt{max\_num\_seqs}=64, and 16-token blocks.
The second term of Eq.~\ref{eq:appendix:frontier-depth} is deployment-specific
and is recomputed from the calibrated service profile.

The frontier is not a protected queue for predicted SLO misses.
If a dispatch-eligible request fits the capacity budget, it proceeds;
otherwise it remains Router-owned, prospectively accounted, and eligible for
replanning.
Rejected registration follows the handoff rollback semantics in
Section~\ref{sec:appendix:handoff}.


% --------------------------------------------------------------------------
\subsection{Engine-Local TBT Protection}
\label{sec:appendix:tbt-control}

RAVEL's engine adapter protects the local execution assumptions used by
LATENCY quotes.
The adapter is installed only when the target vLLM exposes the expected
scheduler hooks; it fails fast if
\texttt{Scheduler.\_get\_priority} or
\texttt{Scheduler.\_schedule\_chunked\_prefill} is unavailable rather than
silently running an unvalidated compatibility path.

Let
\begin{equation}
    T_r
    =
    \min_{i\in\mathcal{D}^{\mathrm{lat}}_r}
       D_i^{\mathrm{TBT}}
    \label{eq:appendix:tightest-tbt}
\end{equation}
be the tightest active TBT requirement among LATENCY Decode requests at
replica $r$.
From the calibrated profile, the adapter forms the conservative envelope
\begin{equation}
    a^{\mathrm{guard}}_r=\max_n a_{c(r),n},
    \qquad
    b^{\mathrm{guard}}_r=\max_n b_{c(r),n}.
    \label{eq:appendix:prefill-guard-envelope}
\end{equation}

For block size $g$ and ordinary token limit $B_{\max}$, it computes
\begin{align}
    q_{\mathrm{raw}}
    &=
      \left\lfloor
         \frac{\max(0,T_r-b^{\mathrm{guard}}_r)}
              {a^{\mathrm{guard}}_r}
      \right\rfloor,
      \label{eq:appendix:chunk-raw}\\
    q_{\mathrm{blk}}
    &=
      \min\!\left(
        B_{\max},
        \max\!\left(
          g,\;
          g\left\lfloor\frac{q_{\mathrm{raw}}}{g}\right\rfloor
        \right)
      \right).
      \label{eq:appendix:chunk-aligned}
\end{align}

If $n_r^{\mathrm{dec}}$ Decode sequences are active, each contributes one
next-token Decode slot to the next scheduling iteration.
The temporary block-aligned batch budget is therefore
\begin{equation}
    B^{\mathrm{guard}}_r
    =
    g
    \left\lceil
       \frac{n_r^{\mathrm{dec}}+q_{\mathrm{blk}}}{g}
    \right\rceil .
    \label{eq:appendix:guard-batch-budget}
\end{equation}
The adapter applies
\begin{equation}
    B_{\mathrm{iter}}
    \leftarrow
    \min(B_{\max},B^{\mathrm{guard}}_r)
    \label{eq:appendix:temporary-batch-cap}
\end{equation}
for that Chunked-Prefill scheduling call only and restores $B_{\max}$
afterward.

If no active Decode request carries a LATENCY TBT requirement, the adapter
uses the ordinary vLLM scheduling path.
If even one minimum-size block exceeds the profiled TBT budget, the adapter
must still expose a legal vLLM block and therefore cannot claim a TBT
guarantee in that physical regime.
The Router feasibility test is intended to avoid such placements, and any
realized violation remains in the evaluation denominator.

Phase-aware priority metadata carries a remaining deadline budget rather than
a Router monotonic-clock timestamp.
The engine reconstructs its local deadline using its own arrival clock,
avoiding an implicit cross-host monotonic-clock synchronization assumption.
Non-RAVEL priorities occupy a separate range, so the adapter remains inert for
ordinary vLLM work.


% --------------------------------------------------------------------------
\subsection{Commitment and Engine-Handoff State Machine}
\label{sec:appendix:handoff}

RAVEL defines commitment as a successful ownership transition rather than a
timer expiration, queue removal, or registration attempt.
A request may be tentatively placed, held, replanned, and made
dispatch-eligible while remaining Router-owned.

The engine-facing handoff performs four operations:
(1) freeze the latest Router quote for observability;
(2) ask the tentative target to register the request;
(3) create the asynchronous engine-facing posting task; and
(4) install that task in the attempt-indexed task table.
Quote freezing is diagnostic and does not commit placement.
Successful registration together with successful creation and tracking of the
posting task is the sole commitment event.

\begin{table*}[t]
    \centering
    \small
    \caption{Engine-handoff state transitions.}
    \label{tab:appendix:handoff-state-machine}
    \begin{tabular}{
        p{0.16\textwidth}
        p{0.20\textwidth}
        p{0.43\textwidth}
        p{0.10\textwidth}}
        \toprule
        \textbf{Phase} &
        \textbf{Event} &
        \textbf{Action} &
        \textbf{Placement final?} \\
        \midrule

        Router-owned &
        Dispatch claim &
        Enter \texttt{COMMITTING} and attempt target registration. &
        No \\

        \texttt{COMMITTING} &
        Registration rejected &
        Restore Router ownership and prospective accounting. &
        No \\

        \texttt{COMMITTING} &
        Posting setup fails &
        Abort partial registration and clean attempt-specific state. &
        No \\

        \texttt{COMMITTING} &
        Registration and posting-task tracking succeed &
        Transfer ownership to the serving engine. &
        Yes \\

        Engine-owned &
        Completion, failure, or cancellation &
        Perform attempt-specific terminal cleanup; do not reopen placement. &
        Yes \\

        Any &
        Stale callback &
        Ignore the callback when its attempt identifier is not current. &
        Unchanged \\

        \bottomrule
    \end{tabular}
\end{table*}

Before successful handoff, failure may restore Router-owned state or clean a
partial attempt without reporting a committed placement.
After successful handoff, failures belong to the committed attempt and do not
create another geographic-placement opportunity.

Task cleanup is attempt-aware.
A deletion or terminal callback carrying attempt identifier $h$ is ignored
when the stored task belongs to a different active attempt.
This prevents delayed callbacks from releasing ownership or accounting state
associated with a newer attempt.

Before successful handoff, RAVEL changes only Router/transport metadata and
prospective accounting; engine-managed execution state has not yet been
materialized in the form that would require runtime KV-state migration.
After handoff, pending and running engine state remains visible to subsequent
quotes, but the committed request itself is no longer an ordinary geographic
replanning candidate.


% ==========================================================================
% Appendix B
% ==========================================================================
